\section{Declaration of the use of AI and outline of the process}
\label{app:ai-declaration}

This article was prepared autonomously by AI agents, as instructed by Gustavo
Jasso, between 7 and 8 October 2026. All its text, proofs, computations and
formal proofs were produced by AI models; no part of it was written by a human.
Gustavo Jasso designed the task and took the decisions recorded below. No human
has checked the mathematics of this article, and it is not peer reviewed.

\subsection*{Models}

The models and their roles were as follows; the names are those of the
project's records.
\begin{itemize}
  \item Claude Opus~5.5 (Anthropic), run in Claude Code: coordination of the
  work, the analysis of the preprints, the general criteria and obstructions
    of \Cref{sec:criteria} and \Cref{sec:obstructions} (part of them jointly
    with the other models, as recorded in the repository), the proof notes
    compiled from earlier results, the checking of verification reports, the
    first stage of the Lean formalisation, and the first draft of the prose of
    this article.
  \item GPT-6 Astra (OpenAI), run in Codex (command-line interface, version
    0.161), with reasoning effort between medium and its highest setting:
    independent verification of every proof in fresh sessions, all heavy
    computations, the proof notes for the constructions of the two preprints,
    the Tate-duality obstruction of \Cref{thm:tate-obstruction}, the later
    stages of the Lean formalisation, and the block-by-block review of the
    article.
  \item Claude Fable~5.1 (Anthropic): two consultations on possible simpler
    constructions and on the obstructions to them; the second was incomplete.
  \item Claude Sonnet (Anthropic), as subagents: retrieval of literature and
    bibliographic data, summaries of instructions, and one computation; their
    output was checked before use. In the second checks described in
    \Cref{app:verification}, Claude Sonnet~5.5 coordinated the work and
    drafted several edits, which were reviewed by Claude Opus~5.5.
\end{itemize}
A ChatGPT Pro~100 subscription and a premium seat in the Claude Team plan for
scientists (at a discounted rate) were available for the work.
% Usage figures: USAGE.md in the repository.

\subsection*{Process}

The first part (7--8 October 2026) analysed the main
preprint~\cite{OAI26findim} and searched for a simpler counterexample to the
little finitistic dimension conjecture; it produced the criteria of
\Cref{sec:criteria} and the obstructions of \Cref{sec:obstructions}, but found
no counterexample verifiable by hand. The second part
produced this article from that record.

Every mathematical claim carried a status (heuristic, plausible, supported by
computation, proved by an AI model, or verified by a second, independent check
of a different kind) and every change to a stronger status was accompanied by a
record of the supporting evidence. Proofs were checked step by step in fresh
sessions, sometimes of
the same model that wrote them. For the checks of the proof notes, earlier
verdicts were removed as far as possible, and the reviewers were instructed not
to use any remaining assessment as evidence; several errors and gaps were found
and
corrected in this way, and two broader claims were withdrawn. Computations used
exact arithmetic over finite fields and rational function fields, and their
scope is stated wherever they are used (\Cref{app:computations}). Citations were
checked against the sources, whose statements and locators were read before
being cited. Parts of \Cref{sec:criteria,sec:selection,sec:realisation} are
formally verified in Lean against Mathlib, and the arguments of
\Cref{sec:obstructions} are formally verified under explicitly stated
hypotheses (\Cref{app:verification}). The completed article was reviewed block by
block by fresh sessions of GPT-6 Astra without waiting for human approval of
each block, in three parts. The reviewers checked
$337$ blocks, corrected punctuation, typesetting and a few slips directly, and
proposed $15$ further changes, concerning notation and the precision of
statements about scope, attribution and evidence; each proposal was examined
by Claude Opus~5.5, and all were adopted. No error in a proof was found at this
stage.

\subsection*{Repository}

The repository
\url{https://github.com/gjassoah/findim-report} contains the complete record of
the work: the plan, the
notes, the verification reports, the scripts and their outputs, the Lean
sources and the evidence of their checks, the sources of this article, and a
log of all interactions between Gustavo Jasso and the agents, in which his
messages are reproduced verbatim and the agents' replies are summarised.
